\section{Measurement discrimination, geometry, and learning}
\label{sec:matrixcomparison76}
The present binary interface retains the ordered classical outcome
and an arbitrary external reference, but no postmeasurement quantum
system from the device. This fixes the comparison with the
measurement-discrimination literature. The projective theorem of
Pucha\l a and collaborators \cite[Corollary~1 and Theorem~2]{PPKKO72}
gives an exact multiple-use formula with ancillary access. It is
used in the retained observable-readout theorem. The restricted
ancillary interface in Fiur\'a\v sek--Mi\v cuda \cite{FM75} leads
to a different comparison of adaptive and parallel probes.

Krawiec--Pawela--Pucha\l a \cite[Section~4]{KPP76} construct a pair
of nine-outcome qutrit rank-one POVMs that admit perfect adaptive
discrimination with two calls and no perfect discrimination by
any finite parallel experiment. Datta--Biswas--Saha--Augusiak
\cite[Theorems~3--5]{DBSA76} study entanglement-assisted
single-use discrimination, including perfect constructions and a
qubit entanglement advantage. Their qubit construction has three
outcomes. These are direct antecedents for general POVM
discrimination with a retained reference. They do not supply a
full binary-effect family covering law. Our comparison
\eqref{eq:matrixadaptivity76} permits an adaptive advantage and
controls its size uniformly in $N$ for fixed input dimension.

The differential argument in Lemma~\ref{lem:matrixpath76} belongs
to the established horizontal-gauge and channel-extension method
\cite{FI76,DKG67,DM76,KGAD76}. Finite adaptive discrimination
bounds obtained by integrating a local channel bound appear in
\cite{YF76,SD76}. The binary Sylvester choice here gives the
specific regularized variance $M(I-M)$, which can be integrated
explicitly and matched by compressed nonadaptive tests. The
inverse-Sylvester expressions of Bures geometry
\cite{Dittmann76} and its spectral integration
\cite{SZGeom76} are related methods. Our form acts on the effect
displacement $H$ with variance $V=M(I-M)$; it is not the ordinary
Bures pullback under $M\mapsto V$, whose differential would be
$H-MH-HM$. The operational-ball lemma and the nested endpoint
integral are both required for Theorem~\ref{thm:matrixcover76}.

Mele--Bittel \cite[Lemmas III.7--III.8 and Corollary III.9]{MB67}
give the binary one-use identity and an operator-norm learner.
Zambrano--Ramos-Calderer--Kueng \cite{ZRK76} study one-use outcome
variation under their single-copy probe model. Combining a
one-use $\epsilon$ guarantee with the hybrid bound gives a
future-$N$ guarantee after setting $\epsilon$ of order
$\delta/N$. Theorem~\ref{thm:commonlearn76} instead proves and
matches the order $N\delta^{-2}\log(1/\eta)$ for the future
adaptive loss on all of $\mathfrak E_2$. The learning network of
\cite{BDPS76} stores quantum resources for later measurement
retrieval and has a different output interface.

Multiscale phase estimation and robust phase recovery are
established techniques \cite{HB76,KLY76}. Our proof supplies its
own small-cone argument and conditional confidence allocation;
it does not invoke the numerical unwrapping threshold corrected
by the erratum to \cite{KLY76}. Randomized signs remove the
unknown bias exactly, and empirical amplitudes choose the
entangled block length without a supplied noise parameter.
The learning theorem is a common experiment for all unknown
qubit effects. The pair-dependent tests used in the matrix
metric comparison and the supplied-target exact encoder retain
their separate roles.

All matrix entropy constants are for fixed dimension and small
error. The proof includes every rank and spectral multiplicity
of the binary effect body. Multi-outcome POVMs, residual quantum
outputs, growing-dimension constants, and full submaximal-error
matrix coverings require additional arguments. The inherited
full-error preparation and fixed-readout results retain their
own stated ranges.
